\section{Introduction}

\subsection{Background and Motivation}
Retail supply chain management operates in an increasingly complex, data-intensive, and volatile environment where accurate demand estimation is fundamental to operational viability \cite{tiwari2025agentic}. Modern multi-channel retail enterprises must coordinate procurement, warehousing, and inventory allocation across hundreds of stores and thousands of stock-keeping units (SKUs). These retail portfolios exhibit substantial demand heterogeneity, characterized by pronounced day-of-week seasonality, periodic promotional surges, holiday spikes, and intermittent low-velocity purchasing patterns \cite{mitra2022comparative, venkiteela2025machine}. Within this operational setting, demand forecasting serves as the cornerstone of Sales and Operations Planning (S\&OP), directly governing procurement schedules, labor allocation, and working capital commitments \cite{mustapha2025forecasting}. However, demand forecasting is inherently subject to error. Inaccurate demand projections propagate downstream into operational failures: under-forecasting leads to inventory stockouts, immediate lost sales, and customer defection, whereas over-forecasting results in excessive inventory holding costs, working capital lockup, and product obsolescence \cite{spanos2022application, malik2024sales}. While predictive accuracy is a necessary prerequisite for effective planning, operational decisions depend not only on expected demand quantities but also fundamentally on the uncertainty surrounding those predictions.

In conventional enterprise forecasting workflows, predictive models output a single deterministic point forecast, denoted as $\hat{y}_{t+h}$ for a future horizon $h$ from forecast origin $t$. While point forecasts provide an estimate of conditional expected demand, they do not communicate the dispersion, skewness, or tail risk inherent in the forecast error distribution \cite{gneiting2007strictly}. Standard machine learning loss functions, such as Mean Absolute Error (MAE) and Root Mean Squared Error (RMSE), train models under symmetric error penalization, implicitly treating positive and negative forecast errors as equally costly. Yet, in retail replenishment, the financial and operational consequences of estimation errors are fundamentally asymmetric \cite{ban2019big, bertsimas2020predictive}. An under-forecast of ten units during a peak sales window directly triggers stockouts and permanently lost revenue under lost-sales dynamics, whereas an over-forecast of ten units merely incurs a marginal daily holding cost. Consequently, relying exclusively on deterministic point forecasts provides limited information for inventory replenishment decisions that must actively balance asymmetric operational penalties.

To capture the complex, non-linear dependencies governing retail sales, the forecasting literature has progressed rapidly from univariate linear statistical models, such as Autoregressive Integrated Moving Average (ARIMA) and Exponential Smoothing (ETS), to global machine learning algorithms \cite{mustapha2025forecasting, yeesin2024application}. In extensive empirical benchmarks, including the large-scale M5 competition, Gradient Boosted Decision Trees (GBDTs), most notably Extreme Gradient Boosting (XGBoost) and LightGBM, have demonstrated superior predictive accuracy and computational efficiency compared to classical statistical benchmarks and recurrent neural architectures on structured tabular demand series \cite{makridakis2022m5, chen2016xgboost, hobor2026comparative}. Global tree-based models effectively pool cross-series information across large store-item portfolios, capturing shared seasonal patterns and category-level trends \cite{alghurabli2025ai}. However, in multi-step forecasting settings, standard recursive forecasting strategies, where model predictions are iteratively fed back as lagged predictors, frequently suffer from compounding error propagation and unintended look-ahead feature leakage if feature pipelines access observations within the forecast lead-time window \cite{tiwari2026explainable}. Ensuring that multi-horizon feature pipelines remain strictly decision-time-safe is therefore critical for real-world operational integrity.

\subsection{Conference Baseline and Research Problem}
In our prior conference baseline study \cite{e_idss_ieee}, we introduced the Explainable Intelligent Decision Support System (E-IDSS), which deployed an optimized single-step XGBoost forecasting model integrated with TreeSHAP (Tree-based Shapley Additive exPlanations) \cite{lundberg2017unified} to provide post-hoc local and global feature attributions for retail sales predictions. The E-IDSS framework addressed the enterprise explanation gap, demonstrating how transparent feature importances can demystify black-box machine learning models for store managers and reduce subjective forecast override rates \cite{lin2026critical, narlawar2025explainable, chang2025explainable}. Nevertheless, a critical methodological boundary remained: post-hoc explainability describes why a historical forecast was generated by attributing importance to past input features, but it does not by itself quantify forward predictive uncertainty or determine how forecast risk should be translated into physical replenishment order quantities. Bridging predictive machine learning with prescriptive inventory decision support requires complementing feature-level transparency with rigorous, calibrated uncertainty bounds.

To quantify forecast dispersion without imposing restrictive parametric assumptions, research in probabilistic forecasting and distribution-free uncertainty quantification has increasingly turned to Conformal Prediction (CP) \cite{vovk2005algorithmic, angelopoulos2021gentle}. Unlike traditional safety stock formulations that assume normally distributed, identically and independently distributed (i.i.d.) forecast errors \cite{silver1998inventory}, conformal prediction provides a mathematically rigorous, model-agnostic framework for constructing valid prediction intervals using empirical nonconformity scores derived from historical calibration residuals \cite{stankeviciute2021conformal}. When applied to time series, rolling-origin conformal calibration pools allow non-parametric quantile estimation that adapts to observed residual dispersions \cite{huber2019integrating}. However, standard conformal prediction theory fundamentally relies on the assumption of data exchangeability, an assumption that is frequently violated in nonstationary retail environments \cite{barber2023conformal}.

Real-world retail demand is notoriously nonstationary, exhibiting macroeconomic trend shifts, evolving consumer purchasing habits, and annual scale expansions \cite{malik2024sales}. In the presence of forward temporal distribution shifts, the statistical dispersion of future forecast errors can diverge substantially from historical calibration residuals \cite{tibshirani2019conformal}. Under such conditions, standard rolling conformal prediction intervals may experience empirical coverage degradation, failing to achieve nominal target coverage levels (such as 80\%, 90\%, or 95\%). Rather than assuming away nonstationarity, it is essential to empirically characterize the extent of coverage degradation under temporal shift and evaluate whether scale-adaptive nonconformity scores can maintain operational utility in downstream decision systems.

The ultimate value of demand forecasting is realized not in statistical isolation, but through its integration into inventory replenishment decisions \cite{syntetos2016supply}. In classical continuous-review $(r, Q)$ and $(s, S)$ inventory control, replenishment orders are triggered whenever the inventory position (on-hand plus on-order stock) drops below a Reorder Point ($ROP$) \cite{axsaeter2015inventory}. Over a supplier replenishment lead time $L$, the reorder point must cover both the expected cumulative lead-time demand $\hat{D}_{t,L}$ and a safety stock buffer $SS_{t,L}$ to protect against stockout risk: $ROP_{t,L} = \hat{D}_{t,L} + SS_{t,L}$. While extensive operations research literature evaluates inventory policies under theoretical demand distributions \cite{snyder2019fundamentals}, and modern machine learning literature evaluates prediction intervals using scoring rules such as the Winkler score \cite{gneiting2007strictly}, the operational bridge connecting non-parametric conformal uncertainty margins directly to dynamic safety stock buffers across multi-day lead times and lost-sales dynamics remains underexplored. The critical research question is not merely whether a prediction interval achieves statistical calibration, but how the resulting uncertainty-derived replenishment buffers perform when subjected to continuous discrete-event operational simulation.

\subsection{The Research Gap}
A critical synthesis of the literature reveals four interconnected research gaps spanning forecasting, uncertainty quantification, and inventory control:

\begin{enumerate}
    \item \textbf{Gap 1: Forecast Accuracy versus Downstream Decision Quality}: Existing retail forecasting studies predominantly prioritize predictive error minimization (MAE, RMSE, WAPE) on static benchmarks \cite{makridakis2022m5, hobor2026comparative}. However, lower point-forecast errors do not inherently guarantee superior inventory decisions, as symmetric error metrics do not reflect the asymmetric financial penalties of lost sales versus holding costs \cite{ban2019big, bertsimas2020predictive}. Limited empirical evidence connects strictly leakage-free multi-horizon retail forecasting with downstream replenishment outcomes using a common decision-level evaluation framework.
    \item \textbf{Gap 2: Conformal Uncertainty Calibration Under Temporal Distribution Shift}: While conformal prediction provides distribution-free finite-sample guarantees under exchangeability \cite{vovk2005algorithmic, angelopoulos2021gentle}, retail sales series routinely exhibit temporal nonstationarity and variance expansion \cite{barber2023conformal}. Limited evidence exists on how forecast-conditional uncertainty calibration behaves under forward distribution shift and how empirical uncertainty margins should be structured and interpreted when empirical coverage deviates from nominal confidence targets.
    \item \textbf{Gap 3: Operational Translation of Forecast Uncertainty into Replenishment Buffers}: Existing literature frequently isolates probabilistic forecasting from inventory execution \cite{stankeviciute2021conformal, huber2019integrating}. The operational translation of multi-horizon forecast uncertainty into dynamic safety stock buffers remains insufficiently characterized, particularly regarding how uncertainty margins scale across multi-day lead times ($L \in \{1, 3, 5, 7\}$ days) in lost-sales retail environments.
    \item \textbf{Gap 4: Portfolio-Level Cost-Service Efficient Frontier Characterization}: Inventory optimization involves conflicting objectives: maximizing customer service and eliminating stockouts while minimizing holding costs and order placement overhead \cite{silver1998inventory}. There remains limited portfolio-level empirical evidence showing how uncertainty-aware replenishment affects customer fill rate, lost sales volume, average inventory investment, and total operational costs across varying lead times and shortage penalty regimes.
\end{enumerate}

Synthesizing these dimensions, the central research gap addressed by this study is formulated as follows:

\begin{quote}
\textit{Existing research has substantially advanced machine-learning demand forecasting, explainable artificial intelligence, conformal uncertainty quantification, and inventory replenishment optimization; however, these research streams remain only partially connected in retail decision-support research. Forecasting studies commonly emphasize point predictive accuracy, uncertainty studies emphasize interval calibration under exchangeability, and inventory studies evaluate replenishment policies under stylized parametric demand assumptions. Consequently, there is limited empirical evidence on how leakage-free multi-horizon retail forecasts and forecast-conditional uncertainty estimates behave under forward distribution shift, and how this uncertainty translates into operational replenishment performance across different lead times and shortage-cost regimes. In particular, the empirical relationship among conformal uncertainty calibration, customer fill rate, lost sales prevention, inventory investment, and total simulated cost remains insufficiently characterized at the store-item portfolio level.}
\end{quote}

\subsection{Substantial Extension from Conference Work}
This journal manuscript represents a major methodological, mathematical, and empirical extension of our preliminary conference baseline \cite{e_idss_ieee}. While the conference paper established an initial explainable forecasting architecture using single-step XGBoost, TreeSHAP, and heuristic Multi-Criteria Decision Analysis (MCDA) scoring, this journal extension expands the scope into a fully realized, leakage-free, conformal uncertainty-driven, and discrete-event inventory decision-support framework.

The structural extensions introduced in this manuscript compared to the conference baseline are summarized in Table~\ref{tab:conference_vs_journal}.

\begin{table}[htbp]
\centering
\caption{Methodological and Empirical Extension from Conference Baseline to Journal Study}
\label{tab:conference_vs_journal}
\small
\resizebox{\linewidth}{!}{%
\begin{tabular}{lll}
\toprule
\textbf{Dimension} & \textbf{Conference Baseline (E-IDSS) \cite{e_idss_ieee}} & \textbf{Journal Extension (This Study)} \\
\midrule
Primary Focus & Point forecasting \& heuristic priority ranking & End-to-end forecast uncertainty to inventory decisions \\
Forecasting Architecture & Single-step global XGBoost regression model & 7 Decoupled Global Direct Multi-Horizon XGBoost models ($M_1..M_7$) \\
Feature Temporal Protocol & Standard time-series lag and rolling indicators & Strictly decision-time-safe protocol (minimum lag $\ge 7$ days) \\
Calibration Pool & In-sample residuals / heuristic variance & Rolling Out-of-Sample (OOS) calibration pool ($N=608,165$) \\
Uncertainty Quantification & Heuristic MCDA Risk Score (volatility index) & Rolling Forecast-Conditional Conformal Calibration ($W=100, K^*=100$) \\
Distribution Shift Analysis & Assumed stationary operational environment & Location-scale \& violation-symmetry decomposition of test shift \\
Robustness Analysis & Fixed single heuristic scoring rule & Full calibration neighborhood sensitivity sweep ($K \in [10, 100]$) \\
Inventory Decision Model & Categorical recommendations (Increase/Maintain/Cut) & Quantitative continuous-review order-up-to policy ($ROP = \hat{D} + SS$) \\
Simulation Engine & No physical inventory simulation executed & 8-step daily discrete-event simulation with Lost Sales dynamics \\
Lead-Time Evaluation & Undefined / single-period implicit assumption & Evaluated across 4 multi-day lead times ($L \in \{1, 3, 5, 7\}$ days) \\
Financial Cost Accounting & No cost accounting or financial settlement & Comprehensive cost model ($K_{\text{order}}=\$10, h=\$0.20/\text{day}, p=\$1.00$) \\
Shortage Cost Sensitivity & Fixed heuristic assumption & Shortage penalty sensitivity grid ($p \in \{\$0.60, \$1.00, \$2.00, \$4.00\}$) \\
Statistical Evaluation & Point error metrics (MAE, RMSE, $R^2$) & Paired Wilcoxon signed-rank ($p < 10^{-75}$) + 1,000 bootstrap CIs \\
Evaluation Scale & Static portfolio summary & 500 store-item series simulated over 365 out-of-sample days \\
\bottomrule
\end{tabular}%
}
\end{table}

As detailed in Table~\ref{tab:conference_vs_journal}, the journal extension expands the preliminary baseline across fourteen critical operational and empirical dimensions.

\subsection{Research Objectives and Questions}
The overarching research objective of this study is:

\begin{quote}
\textit{To develop, empirically validate, and economically evaluate an end-to-end decision-support framework that bridges leakage-free direct multi-horizon gradient boosted forecasting, rolling forecast-conditional conformal uncertainty quantification, and discrete-event inventory replenishment simulation across 500 heterogeneous retail store-item series under forward distribution shift.}
\end{quote}

To achieve this objective, the study investigates four specific Research Questions (RQs):

\begin{itemize}
    \item \textbf{RQ1 (Forecasting Architecture and Decision-Time Causality)}: \textit{How does a decision-time-safe, direct multi-horizon gradient boosting architecture perform across varying forecast horizons ($h=1,\dots,7$) in a high-dimensional retail demand environment, and how does the exclusion of look-ahead lag features affect predictive accuracy?}
    \item \textbf{RQ2 (Conformal Calibration Under Distribution Shift)}: \textit{How does rolling forecast-conditional conformal uncertainty calibration behave under forward retail demand distribution shift, and why does empirical coverage deviate from nominal confidence targets in nonstationary time series?}
    \item \textbf{RQ3 (Operational Service and Shortage Performance)}: \textit{How does incorporating non-parametric conformal uncertainty margins into replenishment reorder points affect customer fill rate, stockout frequency, lost sales volume, and average inventory investment across varying lead times ($L \in \{1, 3, 5, 7\}$) under lost-sales dynamics?}
    \item \textbf{RQ4 (Economic Trade-Offs and Shortage Cost Sensitivity)}: \textit{How does the total operational cost of uncertainty-aware replenishment compare to deterministic point-forecast replenishment across different lead-time constraints and asymmetric shortage penalty regimes ($p \in \{\$0.60, \$1.00, \$2.00, \$4.00\}$)?}
\end{itemize}

\subsection{Summary of Primary Contributions}
This manuscript provides five primary scientific, methodological, and empirical contributions:

\begin{enumerate}
    \item \textbf{Decision-Time-Safe Multi-Horizon Forecasting Protocol}: We develop and validate a direct multi-horizon forecasting architecture utilizing seven decoupled global XGBoost regressors ($M_1,\dots,M_7$) equipped with a feature-generation protocol that enforces a minimum lag constraint of $\ge 7$ days, eliminating recursive error accumulation and guaranteeing that zero future target sales enter forecast feature matrices across all evaluation horizons.
    \item \textbf{Rolling Forecast-Conditional Conformal Uncertainty Layer}: We implement a distribution-free conformal calibration layer that dynamically scales nonconformity scores by point-forecast magnitudes ($s_{r,L} = |y_{r,L} - \hat{D}_{r,L}| / (\hat{D}_{r,L} + \epsilon)$) derived from an out-of-sample calibration pool, significantly reducing mean Winkler interval loss ($p < 0.02$) and boosting high-demand coverage by $+8.80$ percentage points compared to unscaled conformal baselines.
    \item \textbf{Empirical Characterization of Conformal Shift Dynamics and Robustness}: We conduct an extensive location-scale diagnostic decomposition demonstrating that out-of-sample conformal coverage degradation is driven by forward residual-scale expansion ($3.55\times$ to $4.30\times$ increase in standard deviation) and seasonal phase misalignment rather than simple point-forecast bias, and establish that calibration-window performance exhibits a broad stability plateau across neighborhood sizes $K \in [30, 100]$.
    \item \textbf{Operational Translation of Forecast Uncertainty into Replenishment Buffers}: We formulate an operational decision-support mechanism that directly couples forecast-conditional conformal uncertainty margins to continuous-review reorder points ($ROP_{t,L} = \hat{D}_{t,L} + SS_{t,L}$), demonstrating in discrete-event simulation across 500 store-item series that conformal safety stock yields statistically significant fill rate gains ($+7.29$ to $+10.64$ percentage points, paired Wilcoxon $p < 10^{-75}$) and cuts lost sales shortages by $20.03\%$ to $29.57\%$ (saving up to $1.14$ million units of unmet demand).
    \item \textbf{Portfolio-Level Empirical Cost-Service Characterization}: We characterize the economic trade-offs of uncertainty-aware replenishment across four lead times and four shortage penalty levels, demonstrating lower total simulated operating costs at $L=1$ (demand-weighted fill rate $+10.64$ pp, total cost $-10.83\%$, saving $\$790,834$, $p = 6.65 \times 10^{-55}$) and establishing an empirical trade-off frontier at $L \ge 3$ where conformal replenishment becomes cost-preferred whenever shortage penalties exceed $\$2.00/\text{unit}$.
\end{enumerate}

\subsection{Paper Organization}
The remainder of this paper is organized as follows: Section 2 reviews related literature in retail forecasting, explainable artificial intelligence, conformal prediction under distribution shift, and inventory replenishment, establishing the theoretical foundations and gap positioning. Section 3 details the end-to-end methodology, including the direct multi-horizon forecasting architecture, out-of-sample residual generation, rolling forecast-conditional conformal calibration, discrete-event inventory simulation engine, and inferential statistical protocol. Section 4 presents the experimental results and discussion across forecasting validation, conformal uncertainty calibration, inventory service levels, cost breakdowns, shortage penalty sensitivity, statistical significance tests, and conceptual implications. Section 5 discusses managerial takeaways, theoretical implications, and study limitations. Section 6 presents formal study limitations. Section 7 provides the conclusion and future research directions.
